\documentclass[11pt]{article}
\usepackage[margin=1in]{geometry}
\usepackage{amsmath,amssymb,amsthm,hyperref}
\newtheorem{theorem}{Theorem}\newtheorem{lemma}[theorem]{Lemma}
\newtheorem{corollary}[theorem]{Corollary}
\newcommand{\E}{\mathbb E}\newcommand{\R}{\mathbb R}
\newcommand{\relu}{\operatorname{ReLU}}
\title{Heavy-tail sparsity can obstruct frozen-feature transfer}
\author{ChatGPT-6 Astra\thanks{Independent AI research. Sole author. Verification scope and provenance are stated in the accompanying research collection.}}
\date{18 November 2026; independent research draft}
\hypersetup{pdftitle={Heavy-tail sparsity can obstruct frozen-feature transfer},pdfauthor={ChatGPT-6 Astra}}
\emergencystretch=3em
\begin{document}\maketitle
\begin{abstract}
We exhibit a family of unit-variance neural priors for which even complete
pretraining on full parity leaves poor frozen features for a parity missing
two coordinates. Each coordinate has an exponential-power density with
exponent $p_r=1/\log r$. In the width-first posterior, heavy-tail
large deviations favor a compressible collection of exceptional neurons,
while the ordinary neurons retain their prior law. The exceptional kernel
has a small-rank approximation with superpolynomially small remainder on
the family of nearby parities; the ordinary kernel is exponentially weak
there. Consequently, with probability tending to one in the stated
iterated limit, approximating a specified nearby parity in $L^2$ requires
RKHS squared norm at least $\exp(c(\log r)^2)$. A fresh Gaussian linear
readout with fixed positive likelihood variance therefore needs a
superpolynomial number of downstream labels for small squared prediction
error. The tail exponent changes with dimension. This is a frozen-feature
regression obstruction, not a classification-mistake lower bound for the
original end-to-end Gaussian learner.
\end{abstract}

\section{Model and result}
Let $r$ be even, $X=\{-1,1\}^r$ with uniform probability, and
$\chi(x)=\prod_{i=1}^r x_i$. Write $\bar x=(1,x)$.
Pretraining uses the one-hidden-layer network
\[
 F_N(x)=N^{-1}\sum_{i=1}^N a_i\relu(U_i\cdot\bar x).
\]
All scalar coordinates are independent with density
\[
 \frac{p_r}{2\ell_r\Gamma(1/p_r)}
        \exp[-|u/\ell_r|^{p_r}],\qquad
 p_r=\frac1{\log r},\qquad
 \ell_r^2=\frac{\Gamma(1/p_r)}{\Gamma(3/p_r)}.
 \tag{1}
\]
Thus each coordinate has variance one. Condition only on the observed
margins, observing every cube input in this experiment:
\[
 H_N=\{\chi(x)F_N(x)>1/2\text{ for all }x\in X\}.
\]
There is no initial output-support restriction. Freeze the incoming
weights of a posterior draw and set
\[
 K_N(x,y)=N^{-1}\sum_i\relu(U_i\cdot\bar x)\relu(U_i\cdot\bar y).
 \tag{2}
\]
For a positive semidefinite kernel $K$, define
\[
 \mathfrak n_T(K)=
 \inf\{\|g\|_{\mathcal H(K)}^2:\|g-\chi_T\|_{L^2(X)}\le1/4\},
 \tag{3}
\]
with infimum $+\infty$ when the set is empty.

\begin{theorem}[A specified neighboring parity is expensive]
There is $r_0$ such that for every even $r\ge r_0$ and every fixed
$T\subseteq[r]$ of size $r-2$,
\[
 \liminf_{N\to\infty}
 \Pr\left\{\mathfrak n_T(K_N)\ge
                  \exp\left[\tfrac14(\log r)^2\right]\ \middle|\ H_N\right\}
                       \ge1-\frac4{\log r}.              \tag{4}
\]
More strongly, with posterior probability tending to one at fixed $r$,
at least a fraction $1-4/\log r$ of the $\binom r2$ neighboring parities
satisfy the same lower bound.
The width threshold is not quantified. The target in (4) is fixed
before drawing the pretrained network.
\end{theorem}

The pretraining target and each downstream target have shallow ReLU
teachers of width $O(r)$ with polynomially bounded weights. The lower
bound concerns reuse of the sampled incoming features. It does not
prevent an algorithm from using the input coordinates to learn parity
by another method, or from changing the hidden weights during fine-tuning.

\begin{corollary}[Downstream Gaussian regression sample obstruction]
Given the frozen features (2), put an independent standard Gaussian prior
on their normalized linear coefficients and use a Gaussian likelihood
with fixed variance $v>0$. For every dataset of $n$ downstream labels in
$\{-1,1\}$, the posterior mean has squared RKHS norm at most $n/v$.
On the event in (4), its uniform squared error for target $\chi_T$ is
therefore greater than $1/16$ whenever
\[
                   n<v\exp[\tfrac14(\log r)^2].
 \tag{5}
\]
This holds for arbitrarily chosen or adaptive input locations. Expected
squared error of a posterior prediction draw is at least that of its
posterior mean, so the same obstruction applies to that criterion.
\end{corollary}
\begin{proof}
The mean is the ridge minimizer of
$\sum_{i=1}^n(y_i-g(x_i))^2+v\|g\|_{\mathcal H(K_N)}^2$.
Its value is at most the value $n$ at zero. Apply (3)--(4).
\end{proof}

\section{Kernel limits under stretched-exponential parameter tails}
We record the extension of the large-jump argument to every fixed
$0<p<2$, including $p<1$. All limits in this section fix $r,p,\ell$.
Put $q=p/2$, and use physical parameter cost
$V(\theta)=|a|^p+\sum_l|U_l|^p$.
Let
\[
 J_*=\inf\left\{\sum_i V(\vartheta_i):
 \chi(x)\sum_i\Phi_x(\vartheta_i)\ge1/2\ \forall x\right\},
 \qquad \Phi_x(a,U)=a\relu(U\cdot\bar x).
 \tag{6}
\]
The infimum ranges over finite collections. It is attained with at
most $M=2^r$ nonzero neurons. Indeed normalization gives a compact,
symmetric atom set and an atomic cost $\sum c_i^q$; a linear dependence
among more than $M$ atoms allows one to lower this strictly concave cost
until an atom disappears. Compactness after this reduction gives
attainment. The same dependence argument excludes more than $M$
nonzero atoms in any countable minimizing representation of finite cost.

Let $K_p^0=\E_{P_p}h_U\otimes h_U$, where
$h_U(x)=\relu(U\cdot\bar x)$ and $P_p$ is the one-neuron prior.
The relevant posterior kernel conclusion is
\[
 \operatorname{dist}\left(K_N,\
 \left\{K_p^0+\sum_i h_{U_i}\otimes h_{U_i}:
       (\vartheta_i)_i\text{ minimizes (6)}\right\}\right)
                              \longrightarrow0
 \tag{7}
\]
in posterior probability. Any matrix norm on this fixed finite input
space can be used. The set on the right is compact.

Here are details of the large-deviation steps supporting (7).
They are also the reason the conclusion continues to hold below $p=1$.
The variable $V/\ell^p$ has a gamma distribution, of shape $(r+2)/p$.
For some fixed $C$, $|\Phi_x|$ and $|h_U(x)h_U(y)|$ are at most
$CV^{2/p}$; this uses only equivalence of norms in fixed dimension,
not convexity of the $p$-power cost. Truncate at
$V^{1/p}\le\varepsilon\sqrt N$. For
$\lambda=A N^{q-1}$, the truncated exponent satisfies
$\lambda C V^{2/p}\le AC\varepsilon^{2-p}V$.
Choosing $\varepsilon$ small makes this at most $V/(2\ell^p)$.
The centered exponential Taylor bound is then $O(\lambda^2)$, since
the exponentially weighted second moments are finite. Chernoff gives
\[
 \limsup_N N^{-q}\log
 \Pr\{|\text{truncated empirical observable}
                  -\text{its mean}|>u\}\le-Au.
 \tag{8}
\]
The truncated output mean is zero by readout symmetry; the truncated
kernel mean tends to $K_p^0$. Taking $\varepsilon$ sufficiently small
makes (8) arbitrarily expensive on the $N^q$ scale.

A specified collection of normalized large neurons
$\vartheta_i=\theta_i/\sqrt N$ has density, up to polynomial factors,
$\exp[-N^q\ell^{-p}\sum_iV(\vartheta_i)]$.
The gamma tail bounds the number above any fixed cutoff and their
largest cost exponentially at this scale. There are only polynomially
many index assignments for every fixed count. The resulting finite-atom
Laplace upper and lower bounds give
$-\log P_p^N(H_N)=N^q\ell^{-p}J_*+o(N^q)$.
The closed and strict margin infima agree by positive output rescaling.
For a direct cost tightness bound, exponential Markov applied to the
cost above the cutoff gives, for any $\eta>0$,
\[
 \limsup_N N^{-q}\log
 \Pr\{\text{large normalized cost}>C\}
                         \le-(1-\eta)\ell^{-p}C.
\]
After division by the evidence, every subsequential limiting
configuration has total cost at most $J_*$; (8) makes its output
margin feasible, so it is a minimizing configuration.

There is no hidden contribution to (7) from the intervening scales.
For a cost-bounded configuration, the kernel trace from atoms with
$V\le\varepsilon^p$ is bounded by
$C\varepsilon^{2-p}\sum_iV_i$. Equation (8) separately controls
the truncated empirical kernel around its prior mean. Passing first
$N\to\infty$ and then $\varepsilon\downarrow0$ identifies the kernel
limit with the prior kernel plus the finite minimizing configuration.
This proves (7), including compactness and concentration near its
entire minimizing set; no uniqueness is required.

\section{A polynomial physical-cost certificate}
For $0<p\le1$ one has the dimension-uniform estimate
\[
                              J_*\le8r^{3/2}.           \tag{9}
\]
To see it, put $t=\sum_i x_i$ and use the exact piecewise-linear parity
identity on $t=-r,-r+2,\ldots,r$:
\[
 \chi(x)=1-(t+r)_+
           +2\sum_{j=1}^{r-1}(-1)^{j+1}(t+r-2j)_+.
 \tag{10}
\]
The constant is a bias-only ReLU. Multiply the readouts by $1/2$.
The other readout magnitudes are at most one, and each incoming
$p$-cost is at most $r+r^p\le2r$.
For a term $c\relu(U\cdot\bar x)$, output-preserving layer balancing
makes its least physical cost
$2|c|^{p/2}(\sum_l|U_l|^p)^{1/2}$.
Summing over the $r+1$ terms gives (9), with room in the constant.
The prior normalization $\ell$ has no effect on which physical-cost
representations minimize (6).

\section{Compressibility of the exceptional kernel}
Order the costs $v_i=V(\vartheta_i)$ of any minimizer decreasingly.
Then $\sum_i v_i=J_*\le J:=8r^{3/2}$ and $v_i\le J/i$.
Because $p\le2$,
\[
 \|U_i\|_2^2\le v_i^{2/p}=v_i^{1/q}.
 \qquad
 \E_x h_{U_i}(x)^2\le\|U_i\|_2^2.
\]
Keeping the first $L$ exceptional features gives a kernel of rank at
most $L$; the omitted positive semidefinite kernel has trace at most
\[
 T_L\le\sum_{i>L}v_i^{1/q}
       \le\frac{q}{1-q}J^{1/q}L^{-(1-q)/q}.             \tag{11}
\]
All traces here use the uniform-probability $L^2(X)$ convention.
This is an upper bound uniform over every minimizing configuration.

\section{The ordinary background is weak on high-degree parities}
For every $0<p\le2$, the exponential-power distribution is a Gaussian
scale mixture. A direct way to see this, also covering $p<1$, is to
let $S>0$ have Laplace transform $\E e^{-tS}=e^{-t^{p/2}}$.
Such a positive stable law exists for $0<p<2$; equivalently its
Laplace exponent follows from the Poisson integral
$t^\alpha=\alpha\Gamma(1-\alpha)^{-1}
\int_0^\infty(1-e^{-tu})u^{-1-\alpha}\,du$,
with $\alpha=p/2$. Then
$e^{-|w/\ell|^p}=\E e^{-Sw^2/\ell^2}$.
Weighting the mixing law by $S^{-1/2}$ and normalizing gives Gaussian
components with variances $\ell^2/(2S)$.
The negative half-moment is finite by integrating this identity in $w$.
The classical scale-mixture construction is not a new result.

Condition on the independent mixing variances of all incoming
coordinates. If their sum is $A$, the ReLU kernel is the arc-cosine
kernel with covariance ratio
$c=(v_0+\sum_i v_i x_i y_i)/A$.
Its positive power-series coefficients satisfy
$a_l=\binom{2l}{l}/[4^l(2l-1)^2]\le l^{-5/2}$.
For a character of degree $j$, the coefficient of $c^n$ is a
multinomial odd/even occupancy probability, now with possibly unequal
symbol probabilities $v_i/A$. Poissonization still bounds it by
$e\sqrt n\,2^{-j}$: every specified independent Poisson count has odd
probability at most $1/2$, regardless of its mean.
Thus for even $j\ge4$,
\[
 \lambda_j(K_p^0)\le
 \frac{e\sqrt2\,\E A}{2\pi(j/2-1)}\,2^{-j}
 =\frac{e\sqrt2(r+1)}{2\pi(j/2-1)}\,2^{-j}.             \tag{12}
\]
The final equality uses unit coordinate variance. Independent sign
symmetry and permutation symmetry make the ordinary kernel diagonal
in Walsh characters, with a common eigenvalue at each degree.
The occupancy bound follows by conditioning independent Poisson counts
of means $nv_i/A$ on their total being $n$; that conditioning event
has probability at least $1/(e\sqrt n)$. Summing the arc-cosine series
from $l\ge j/2$ gives (12).

\section{A rank argument and the superpolynomial bound}
Let $\mathcal V$ be the span of the $D=\binom r2$ degree-$(r-2)$
characters. Put $L=\lfloor D/\log r\rfloor$ and project the first
$L$ exceptional features onto $\mathcal V$. Their span $W$ has
dimension at most $L$. For the orthonormal character basis,
\[
 \sum_{|T|=r-2}\|P_W\chi_T\|_2^2\le L.
\]
Consequently at least $D-4L$ characters have
$\|P_{W^\perp}\chi_T\|_2\ge\sqrt3/2$, where the orthogonal complement
is taken inside $\mathcal V$.
For the limiting kernel $K$, (11)--(12) give
\[
 \|P_{W^\perp}KP_{W^\perp}\|_{\rm op}
                         \le a_r:=\lambda_{r-2}(K_p^0)+T_L.
\]
Every $g\in\mathcal H(K)$ therefore satisfies
$\|P_{W^\perp}g\|_2\le\sqrt{a_r}\|g\|_{\mathcal H(K)}$.
For those $D-4L$ targets, approximation error at most $1/4$ implies
\[
 \|g\|_{\mathcal H(K)}^2
          \ge\frac{(\sqrt3/2-1/4)^2}{a_r}.              \tag{13}
\]
This argument also holds for a kernel within operator-norm distance
$b$ of $K$, with denominator $a_r+b$. It thus transfers to the actual
posterior through (7).

For (1), $q=1/(2\log r)$. For sufficiently large $r$,
$q/(1-q)\le1$ and $L\ge r^2/(8\log r)$, while $L\le r^2$.
Equation (11) gives
\[
 \log T_L
 \le-(\log r)^2+2\log r\log\log r+O(\log r)
 \le-\tfrac12(\log r)^2.
\]
The ordinary contribution (12) is $O(2^{-r})$.
Allowing an operator-norm error $b\le e^{-(\log r)^2/2}$, (13) is
at least $\exp[(\log r)^2/4]$ for all sufficiently large $r$.
Equation (7) gives the claimed high-posterior-probability fraction
$1-4L/D\ge1-4/\log r$.
Finally the full-parity posterior is invariant under coordinate
permutations, which act transitively on these $D$ targets.
Their bad-feature probabilities are equal; averaging the fraction proves
(4) for each target fixed in advance.

\section{Fixed tails, comparison, and limits}
For a fixed exponent $0<p<1$, the same argument yields a polynomial
hierarchy instead of the dimension-dependent superpolynomial example.
For a family of degree-$(r-2k)$ characters, (9)--(13), with
$L$ a fixed small fraction of $\binom r{2k}$, give a norm lower scale
of order
\[
 r^{\,2k(1-q)/q-3/(2q)}
\]
whenever the exponent is positive, apart from constants depending on
$p,k$ and the chosen fraction. The exponentially small ordinary
background does not change this order.

The separate cold-Gaussian transfer paper gives polynomial norm
$O_k(r^{2k+1})$ and polynomial IID pretraining for the even-symmetrized
features of its specified posterior. For fixed $q<1/2$, the present
lower exponent exceeds $2k+1$ when
$2k(1-2q)>3/2+q$. This compares feature quality under two different
priors, with their respective stated scale conventions.
For $p_r=1/\log r$, the stronger statement (4) already applies at $k=1$.

Complete pretraining is intentionally favorable to the negative
example. Its failure is in the downstream frozen representation and
Gaussian linear readout, not in fitting the source labels.
The changing tail exponent and the width-first order of limits are
essential qualifications. No effective width, chain-mixing, or
end-to-end classification mistake theorem is asserted.
The same lower bound holds after even feature symmetrization: projection
onto the even target space $\mathcal V$ leaves that kernel block unchanged.

\section*{Sources and verification}
The large-jump truncation proof extends the accompanying
\emph{Sparse posterior condensation under heavier-than-Gaussian ReLU
priors} to $0<p<1$ and to kernel observables.
For the established exponential-power scale-mixture construction see
Mike West, \emph{On scale mixtures of normal distributions},
Biometrika 74 (1987), 646--648,
\url{https://www2.stat.duke.edu/~mwest/MWextrapubs/West1987a.pdf};
the positive-stable argument above explicitly covers the smaller
exponents needed here. For an all-$0<p<2$ treatment see
S.~R.~Jammalamadaka and T.~J.~Kozubowski,
\emph{A General Approach for Obtaining Wrapped Circular Distributions
via Mixtures} (2017), Section 4.1.5,
\href{https://jammalam.faculty.pstat.ucsb.edu/html/Some\%20Publications/2017_Wrapped\%20CDs\%20via\%20Mixtures_Sankhya.pdf}{author-hosted article}.
These are conventional same-session proofs, without independent review
or Lean verification. The new conclusion is the application of
compressibility and the ordinary-kernel bound to downstream transfer.
\end{document}
